\subsection*{S15. Empty arrays}

A path that is an array in every document, but where no document ever supplies
an element, so no element type is ever observed. If some document does supply
elements, ordinary inference applies and this scenario does not arise: the
element type comes from wherever it was seen.

\begin{lstlisting}[style=json]
{ "id": 1, "tags": [] }
{ "id": 2, "tags": [] }
\end{lstlisting}

\options
\begin{enumerate}[label=\textbf{(\alph*)}]
  \item \textbf{Reject}, with an error naming the path.
  \item \textbf{A child table with an identifier, a parent reference and an
        index, but no value column.} The array exists and is always empty.
        \selected
  \item \textbf{Drop the field.}
  \item \textbf{An opaque raw-JSON element column}, deferring the element type
        to the consumer.
\end{enumerate}

\decision{Emit the child table with no value column, and follow S14: the table
generates no module, and the parent exposes a single accessor.}

\begin{center}
\begin{tabular}{@{}l@{}}
\textbf{Table root} \\
\hline
\texttt{id} \\
\hline
1 \\
\hline
\end{tabular}
\qquad
\begin{tabular}{@{}lll@{}}
\multicolumn{3}{@{}l}{\textbf{Table tags}} \\
\hline
\texttt{id} & \texttt{root\_id} & \texttt{idx} \\
\hline
\multicolumn{3}{@{}l}{\emph{(no rows)}} \\
\hline
\end{tabular}
\end{center}

\begin{lstlisting}[style=ml]
module Root : sig
  type id
  type t = { id : id }
  val get : db -> id -> t
  val all : db -> t list
  val tags : db -> t -> unit list
end
\end{lstlisting}

\rationale{The same principle as S6: keep the field, and let its generated
form state exactly what is known --- here, that the path is an array and that
its element type was never observed. This is again the least element of the
type lattice, and if a later corpus supplies an element the value column
appears and the element type is inferred, exactly as an always-null scalar
path would gain a type.}

\limitation{The table is more degenerate than S6's always-\texttt{None}
column: it has no value column and, in a corpus of this shape, no rows either.
It nonetheless carries one bit of information, because the presence marker
retained in S13 still distinguishes \texttt{\{"tags": []\}} from \texttt{\{\}}
--- \texttt{Some []} against \texttt{None} --- so the table records that the
field was present even though it never records a value.}
